\documentclass[11pt,a4paper]{article}
\usepackage[T1]{fontenc}
\usepackage[utf8]{inputenc}
\usepackage{lmodern}
\usepackage{amsmath,amssymb,amsthm}
\usepackage[margin=27mm]{geometry}
\usepackage{booktabs}
\usepackage[hidelinks]{hyperref}
\hypersetup{pdftitle={Action Differences, Fluctuations and Operational Bounds},pdfauthor={navstokgap research working notes}}
\newtheorem{proposition}{Proposition}
\newtheorem{remark}{Remark}
\newcommand{\dd}{\mathrm{d}}
\title{Action Differences, Fluctuations and Operational Bounds\\
\large Control results for a research programme on mechanical gaps}
\author{navstokgap research working notes}
\date{5 September 2026 --- working draft, not externally peer reviewed}
\begin{document}
\maketitle
\begin{abstract}
We distinguish a nonzero action scale from a gap in action values, an uncertainty
bound and a spectral gap. A constant-force trajectory gives an exact relation
between a geometric error area and a fixed-endpoint action difference. The same
model excludes an intrinsic positive action-value gap on the unrestricted path
space. We give an elementary extension to nondegenerate variations, distinguish
Jacobi fields from arbitrary path fluctuations, and test a proposed delta-derivative
interpretation using the normalized free quantum kernel. Under explicitly quantum
measurement assumptions we derive a finite-resource distinguishability threshold,
which is not a universal gap. These are control calculations, not a derivation of
quantum mechanics or a claim of novelty. They specify what an extended model would
need to change and what remains to be established.
\end{abstract}

\section{Question, notation and evidence status}
The motivating question is whether limitations on Newtonian trajectories can
explain a nonzero action scale, and whether a useful toy model can isolate the
formation of a gap. We write $h=2\pi\hbar$ for Planck's constants, $\eta$ for a path
variation, and $J$ for a second-variation operator. These are distinct objects.
An additional physical field would require its own definition, degrees of freedom
and dynamics; it is not supplied by naming $\eta$ an ``$h$ field.''

For a path class $\mathcal P$ and reference $q_*$, one possible action-gap question is
\[
 g_S=\inf\{ |S[q]-S[q_*]|:q\in\mathcal P,\ S[q]\ne S[q_*]\}>0.
\]
The set must be nonempty and the path class, topology, endpoints and duration
specified. This is different from the Hamiltonian question
$\operatorname{spec}(H)\cap(E_0,E_0+\Delta)=\varnothing$ for some $\Delta>0$.
For the latter, existence of a ground sector, the operator domain and the meaning
of $E_0$ must also be stated. Neither follows merely from an action parameter's units.

Newton's opening geometric constructions motivate our area comparison
\cite{Newton1846,NewtonWilkins}. Our use of Hamilton's action and quantum states is
modern. The repository records reading coverage by edition and passage. It has
not completed a reading of all the Classical Scholia or established that doubts
about vanishing areas delayed the 1687 edition. The newly added miscellaneous
fragment collection NATP00385 \cite{NewtonNATP00385} is not itself a complete
edition of those scholia. Historical motivation is independent of the proofs below.

\section{Constant force and the action-value obstruction}
Fix $m,F,T,v_0>0$, set $V(y)=-Fy$, and impose
$x(0)=y(0)=\dot y(0)=0$, $\dot x(0)=v_0$. Newton's equations give
\[
 x_*(t)=v_0t,\qquad y_*(t)=\frac{F t^2}{2m},\qquad 0\le t\le T.
\]
The kinetic-energy gain, not total-energy change, is
$\delta E=F^2T^2/(2m)$. The triangle between initial tangent, endpoint vertical
and chord has area $v_0FT^3/(4m)$. The chord--curve lens has one third of this:
\[
 \mathcal A_{\rm lens}=\frac{v_0FT^3}{12m}.
\]
These configuration-space areas are frame dependent; they are not canonical
phase-space areas or variances of conjugate observables.

Hold $x=x_*$ fixed and consider the same-endpoint paths $y=y_*+\eta$ with
$\eta\in H^1_0(0,T)$, the Sobolev space with square-integrable weak derivative
and zero endpoint trace. Define
\[
 S[x,y]=\int_0^T\left[\frac m2(\dot x^2+\dot y^2)+Fy\right]\dd t.
\]
The linear variation vanishes by integration by parts and $m\ddot y_*=F$:
\begin{equation}\label{eq:quadratic}
 S[x_*,y_*+\eta]-S[x_*,y_*]=\frac m2\int_0^T\dot\eta^2\,\dd t.
\end{equation}
The straight chord $y_{\rm ch}=FTt/(2m)$ is admissible but is not another solution
under the same constant force. Substitution into \eqref{eq:quadratic} gives
\begin{equation}\label{eq:area}
 \Delta S_{\rm ch,*}=\frac{F^2T^3}{24m}
 =\frac{F}{2v_0}\mathcal A_{\rm lens}=\frac{T\delta E}{12}.
\end{equation}
Adding the same total time derivative $\dd G(q,t)/\dd t$ to both Lagrangians
leaves this difference unchanged because configuration endpoints and duration
coincide. A velocity-dependent endpoint function requires additional endpoint data.

\begin{proposition}[No unrestricted classical action gap]\label{prop:no-gap}
For this fixed-endpoint path class the infimum of positive action differences
from the classical path is zero, although the classical path is the unique
minimizer within the class with fixed horizontal motion.
\end{proposition}
\begin{proof}
If $\eta\ne0$ in $H^1_0$, its derivative cannot vanish almost everywhere, so
\eqref{eq:quadratic} is positive. For $\eta_a(t)=a t(T-t)$ it equals
\[
 \Delta S(a)=\frac{m a^2T^3}{6}.
\]
Any prescribed positive bound is violated by taking a sufficiently small nonzero
$a$. The zero variation alone gives equality in \eqref{eq:quadratic}.
\end{proof}

\begin{proposition}[Local obstruction along a variation]\label{prop:local}
Let $q_a$ be admissible for all $a$ in an open interval about zero, and suppose
$f(a)=S[q_a]$ is twice continuously differentiable there, with
$f'(0)=0$ and $f''(0)=b\ne0$. Then arbitrarily small nonzero absolute action
differences $|S[q_a]-S[q_0]|$ occur. If $b>0$, these differences can be taken positive
without absolute values.
\end{proposition}
\begin{proof}
Taylor's formula gives $f(a)-f(0)=b a^2/2+o(a^2)$. For all sufficiently small
nonzero $a$, its magnitude lies between $|b|a^2/4$ and $3|b|a^2/4$ and its sign
is the sign of $b$. It is nonzero and tends to zero.
\end{proof}
Locality here is initially in the parameter $a$. To conclude that these actions
occur in every path-space neighbourhood of $q_0$, additionally require
$q_a\to q_0$ in the specified path topology. The polynomial family above converges
in $H^1$ (and in $C^1$).

These propositions do not cover discrete path classes, non-smooth restricted
actions, different topological sectors or operational equivalence relations.
Such changes may be interesting, but must be separately motivated. They also do
not resolve existence through singular Kepler collisions: admissible off-shell
paths and solutions of an initial-value problem are different sets.

\section{Jacobi fields and normalized fluctuation gaps}
For $S[q]=\int_0^T(m\dot q^2/2-V(q))\dd t$ about a smooth classical solution,
the second variation is
\[
 Q[\eta]=\int_0^T\left[m\dot\eta^2-V''(q_*)\eta^2\right]\dd t,
 \qquad J=-m\frac{\dd^2}{\dd t^2}-V''(q_*).
\]
For smooth bounded $V''(q_*(t))$ take the Dirichlet operator domain
$H^2(0,T)\cap H^1_0(0,T)$ in $L^2(0,T)$. A variation through classical solutions
satisfies the Jacobi equation $J\eta=0$; arbitrary off-shell variations do not.
For constant force a Jacobi field is $A+Bt$, so fixed zero endpoints force it to
vanish. Equation \eqref{eq:quadratic} nevertheless admits nonzero fluctuations.

For constant force the Dirichlet eigenfunctions are $\sin(n\pi t/T)$ with
eigenvalues $\lambda_n=m(n\pi/T)^2$, $n\ge1$. Thus
\[
 Q[\eta]\ge\frac{m\pi^2}{T^2}\|\eta\|_{L^2}^2.
\]
This is coercivity relative to a norm, not a smallest nonzero value of $Q$.
Rescaling $\eta$ rescales $Q$ quadratically. Moreover $\lambda_n$ has units
mass/time$^2$, while $\|\eta\|_{L^2}^2$ has units length$^2$ time, producing action.
Calling $\lambda_1$ a Planck constant would be dimensionally incorrect.

For the oscillator $V=m\omega^2q^2/2$, the corresponding eigenvalues become
$m[(n\pi/T)^2-\omega^2]$. Zero modes occur at $T=n\pi/\omega$; positivity of
all modes fails after the first such time. This explicit calculation is a useful
next comparison with the quantum oscillator Hamiltonian, whose energy spectrum
requires a separate operator and quantization assumptions. Stationary action
need not mean a global minimum for arbitrary durations and potentials.

\section{What distribution is the oscillatory kernel?}
The quantum path-phase input is $\exp(iS/\hbar)$, with time-integrated action,
as in the Feynman formulation \cite{Feynman1948,FeynmanLectures}. A bare exponential
is not by itself a normalized distributional approximation to a delta derivative.
The free one-dimensional Schr\"odinger kernel provides a precise control case.
For $\hbar,m,\tau>0$ let
\[
 K_\tau(z)=\left(\frac{m}{2\pi i\hbar\tau}\right)^{1/2}
 \exp\left(\frac{imz^2}{2\hbar\tau}\right),
\]
with the usual branch $1/\sqrt{i}=e^{-i\pi/4}$ and oscillatory/tempered-distribution
interpretation. Use Fourier transform
$\widehat f(k)=\int e^{-ikz}f(z)\dd z$. Equivalently this kernel is defined by
\[
 \widehat K_\tau(k)=\exp\left(-\frac{i\hbar\tau k^2}{2m}\right).
\]
For each Schwartz test function, dominated convergence in Fourier space implies
\begin{equation}\label{eq:kernel}
 K_\tau\longrightarrow\delta\quad(\tau\downarrow0,\ \hbar>0\text{ fixed}).
\end{equation}
Taylor expansion of the multiplier, using
$|e^{-iu}-1+iu|\le u^2/2$ for real $u$, further gives
\[
 K_\tau=\delta+\frac{i\hbar\tau}{2m}\delta''+O(\tau^2)
 \quad\hbox{in the sense of pairing with each fixed Schwartz test function.}
\]
The remainder estimate follows by integrating $k^4$ times the absolute value of
the transformed test function. The limit is a delta, and the first correction is
its second derivative, not its first derivative. The free kernel is even in $z$.
One can obtain $\delta'$ by differentiating $K_\tau$ with respect to $z$, but that
is a different distribution. A Fourier constraint integral, for example
$\int e^{i\lambda z/\hbar}\dd\lambda/(2\pi\hbar)=\delta(z)$, also needs an
integration variable and normalization. Substitution of a nonlinear constraint
requires further regularity and Jacobian conditions.

This calculation neither constructs a general real-time path-integral measure nor
proves a stationary-phase theorem on infinite-dimensional path space. It tests
one exact model and keeps the short-time limit separate from semiclassical
reasoning and from any proposed physical cutoff.

\section{A conditional operational action threshold}
Assume the quantum state/measurement framework, a specified positive $\hbar$,
and a coherently controlled two-arm system realizing a relative phase
$\phi=\Delta S/\hbar$. A global phase on a single state is unobservable. Instead
compare two known hypotheses with equal priors:
\[
 |\psi_0\rangle=\frac{|0\rangle+|1\rangle}{\sqrt2},\qquad
 |\psi_\phi\rangle=\frac{|0\rangle+e^{i\phi}|1\rangle}{\sqrt2}.
\]
Allow $N\ge1$ independent identical copies and any joint binary quantum
measurement. This is an idealized hypothesis test, not parameter estimation with
an unknown phase and not a claim about all possible resource protocols.

\begin{proposition}[Finite-copy bound]\label{prop:measurement}
The optimal equal-prior success probability is
\begin{equation}\label{eq:success}
 P_N(\phi)=\frac12\left(1+\sqrt{1-\cos^{2N}(\phi/2)}\right).
\end{equation}
For $1/2<p\le1$, in the first phase lobe $|\phi|\le\pi$, success at least $p$
is possible if and only if
\begin{equation}\label{eq:threshold}
 |\Delta S|\ge d_{N,p}:=
 2\hbar\arccos\left([4p(1-p)]^{1/(2N)}\right).
\end{equation}
For each fixed $1/2<p<1$, $d_{N,p}\to0$ as $N\to\infty$. For $p=1$, the finite-copy
threshold instead remains $\pi\hbar$.
\end{proposition}
\begin{proof}
Two normalized pure states with overlap magnitude $r$ have density-matrix
difference $D$ with eigenvalues $\pm\sqrt{1-r^2}$ on their span (and zero outside).
Indeed $\operatorname{tr}D=0$ and
$\operatorname{tr}D^2=2(1-r^2)$. A binary measurement effect $0\le E\le I$ gives
success $1/2+\operatorname{tr}(ED)/2$, maximized by projecting onto the positive
eigenspace. The maximum is $(1+\sqrt{1-r^2})/2$. Here a single-copy overlap has
magnitude $|\cos(\phi/2)|$, and tensor products raise it to the $N$th power,
proving \eqref{eq:success}. Within the first lobe cosine is nonnegative and
monotone in $|\phi|$, so solving $P_N\ge p$ gives \eqref{eq:threshold}.
For $0<4p(1-p)<1$ its $1/(2N)$ power tends to one. At $p=1$ it is zero for
every finite $N$, giving the stated exception.
\end{proof}

For example $N=1$, $p=3/4$ gives $d_{1,3/4}=\pi\hbar/3$. For fixed
$1/2<p<1$ the asymptotic threshold is
\[
 d_{N,p}\sim2\hbar\sqrt{\frac{-\log(4p(1-p))}{N}}.
\]
There is no monotone lower-bound interpretation over all phases: at $\phi=2\pi n$
the hypotheses coincide again. This model proves a positive operational threshold
at fixed resources and accuracy, not a smallest nonzero action in nature. It
also assumes, rather than derives, the positive $\hbar$ in the phase rule.

If an apparatus realizes the particular relative action \eqref{eq:area}, with
all control and recombination phases included, its first-lobe threshold would be
\[
 T\ge \left(\frac{24m d_{N,p}}{F^2}\right)^{1/3}.
\]
The upper first-lobe condition $F^2T^3/(24m\hbar)\le\pi$ must be retained. The chord
is not an alternative trajectory under the original force alone, so this
substitution is conditional on a physical implementation, not an experimental
prediction for the unmodified projectile problem.

\section{Finite propagation and the field-theory companion}
Finite $c$ alone does not forbid small variations. A concrete control uses the
free relativistic particle, $L=-mc^2\sqrt{1-\dot q^2/c^2}$, fixed endpoints zero,
and $q_a=a t(T-t)$. If $|a|T<c$, the path has a strict speed margin. Relative to
the rest path its action excess is positive for $a\ne0$ and
\[
 \Delta S(a)=mc^2\int_0^T
 \left(1-\sqrt{1-a^2(T-2t)^2/c^2}\right)\dd t
 =\frac{ma^2T^3}{6}+O(a^4)\longrightarrow0.
\]
Thus the speed restriction itself supplies no positive action-value gap in this
class. This example concerns relativistic particle kinematics, not a complete
causal theory of interactions.

Navier--Stokes and Yang--Mills remain complementary comparisons, with their
official targets kept intact \cite{Fefferman,JaffeWitten}. A diffusion or
Dirichlet coercivity bound can depend on time, domain and normalization; a physical
quantum mass gap concerns a different operator and limiting problem. Any future
bridge must identify the spaces, physical versus auxiliary time, estimates and
limits it preserves. Neither the preceding obstruction nor the operational
threshold proves fluid regularity or a four-dimensional quantum-field mass gap.

\section{Reproduction, limitations and next experiments}
The repository's \texttt{make check} verifies the exact polynomial action
identities, the oscillator-mode residual, the Fourier expansion coefficient and
elementary discrimination cross-checks. These tests supplement the displayed
proofs; they are not formal proof certificates. \texttt{make papers} rebuilds this
note and the companion programme PDF. No Lean proof is supplied in this draft.

The next bounded work is to compare oscillator Hessian and Hamiltonian gaps,
free motion on a line and circle, and central-potential collision assumptions.
In parallel, primary-source bibliography will compare quantum reconstruction
axioms and complete the historical corpus inventory. A successful extension must
identify which premise excludes the classical small-amplitude families, or else
choose an operational or spectral gap instead and state the change of target.
Deriving a theory's form, forcing a nonzero universal scale and determining its
measured numerical value are separate tasks.

\bibliographystyle{plain}
\bibliography{references/library}
\end{document}
